% Options for packages loaded elsewhere
\PassOptionsToPackage{unicode}{hyperref}
\PassOptionsToPackage{hyphens}{url}
%
\documentclass[
  10pt,
]{article}
\usepackage{lmodern}
\usepackage{amssymb,amsmath}
\usepackage{ifxetex,ifluatex}
\ifnum 0\ifxetex 1\fi\ifluatex 1\fi=0 % if pdftex
  \usepackage[T1]{fontenc}
  \usepackage[utf8]{inputenc}
  \usepackage{textcomp} % provide euro and other symbols
\else % if luatex or xetex
  \usepackage{unicode-math}
  \defaultfontfeatures{Scale=MatchLowercase}
  \defaultfontfeatures[\rmfamily]{Ligatures=TeX,Scale=1}
\fi
% Use upquote if available, for straight quotes in verbatim environments
\IfFileExists{upquote.sty}{\usepackage{upquote}}{}
\IfFileExists{microtype.sty}{% use microtype if available
  \usepackage[]{microtype}
  \UseMicrotypeSet[protrusion]{basicmath} % disable protrusion for tt fonts
}{}
\makeatletter
\@ifundefined{KOMAClassName}{% if non-KOMA class
  \IfFileExists{parskip.sty}{%
    \usepackage{parskip}
  }{% else
    \setlength{\parindent}{0pt}
    \setlength{\parskip}{6pt plus 2pt minus 1pt}}
}{% if KOMA class
  \KOMAoptions{parskip=half}}
\makeatother
\usepackage{xcolor}
\IfFileExists{xurl.sty}{\usepackage{xurl}}{} % add URL line breaks if available
\IfFileExists{bookmark.sty}{\usepackage{bookmark}}{\usepackage{hyperref}}
\hypersetup{
  hidelinks,
  pdfcreator={LaTeX via pandoc}}
\urlstyle{same} % disable monospaced font for URLs
\usepackage[a4paper,margin=18mm]{geometry}
\setlength{\emergencystretch}{3em} % prevent overfull lines
\providecommand{\tightlist}{%
  \setlength{\itemsep}{0pt}\setlength{\parskip}{0pt}}
\setcounter{secnumdepth}{-\maxdimen} % remove section numbering

\usepackage{esint}
\usepackage{xeCJK}
\setmainfont{DejaVu Serif}
\setmonofont{DejaVu Sans Mono}
\setCJKmainfont{Noto Serif CJK SC}
\setCJKmonofont{Noto Sans Mono CJK SC}
\setmathfont{Asana Math}
\setmathfont{Latin Modern Math}[range=\lrcorner]
\setlength{\emergencystretch}{3em}
\allowdisplaybreaks
\hypersetup{hypertexnames=false}

\author{}
\date{}

\begin{document}

\hypertarget{full-core-cost-signed-extrema-and-conformal-ricci-flow}{%
\section{Full-core cost, signed extrema, and conformal Ricci
flow}\label{full-core-cost-signed-extrema-and-conformal-ricci-flow}}

Research snapshot: 844e62f7304d9f890914a6a4370ec39291ddc26e.

Complete recent proofs. Retain the second paper's full signed target
instead of universally discarding its Weyl payment. Metric terminal-core
extraction does not automatically retain a nonzero actual spinor at that
scale. The flow conclusion is a maximal smooth-time bound, not a
singularity classification.

Read 00\_READ\_FIRST.md for superseded statements. Each source below is
complete and byte-preserved. Original local links remain relative to
their source; locate other texts using Volume 6. Formula numbers restart
in each source.

\begin{center}\rule{0.5\linewidth}{0.5pt}\end{center}

\hypertarget{source-1-20260927_full_core_cost_and_extremal_metrics.md}{%
\subsection{Source 1:
20260927\_full\_core\_cost\_and\_extremal\_metrics.md}\label{source-1-20260927_full_core_cost_and_extremal_metrics.md}}

Original path:
continuation/20260927\_full\_core\_cost\_and\_extremal\_metrics.md

SHA-256:
4625dbfb968571ae91292fc6944d4c969b8a4db681cb8a6cea114ee2529b9302

BEGIN COMPLETE SOURCE

\hypertarget{full-core-curvature-cost-and-extremal-metrics}{%
\section{Full-core curvature cost and extremal
metrics}\label{full-core-curvature-cost-and-extremal-metrics}}

27 September 2026. Continuation from \texttt{0c8059f}.

This proof develops the two open interfaces left in
\href{20260927_local_geometry_and_variational_extrema.md}{the preceding
continuation}: the smooth filling of its actual resonant annuli, and the
all-metric sharp coefficient. It does not reopen the already proved
linear small-h order. No web handoff package is generated.

The filling argument below is a synthesis of classical boundary
curvature formulas with the project's actual holonomy calculation. It is
not a new index theorem or a classification of all Ricci-flat ALE
spaces. Its new use here is that an entire smooth core pays a
scale-independent cost even when its exterior annulus satisfies both
actual local spinor equations.

\hypertarget{an-entire-resonant-core-has-a-uniform-curvature-cost}{%
\subsection{1. An entire resonant core has a uniform curvature
cost}\label{an-entire-resonant-core-has-a-uniform-curvature-cost}}

All orientations in this section are physical orientations, not
orientations chosen after identifying the underlying lens space. Weyl
norms are the Hilbert--Schmidt norms of the endomorphisms of two-forms,
\(|W^\pm|^2=\operatorname{tr}((W^\pm)^2)\), not their spectral operator
norms. Put

\[
 \mathscr Q(g)=|\operatorname{Ric}^{\circ}|^2
                       +4|W^+|^2+\frac{s^2}{12}.
                                                        \tag{1.1}
\]

Let \(\Gamma_n\) act on \(\mathbb C^2\), with its complex orientation,
by \((z_1,z_2)\mapsto(\zeta z_1,\zeta^{-2}z_2)\), where \(n\ge3\) is
odd. Let \(\Omega\) be any compact connected oriented smooth
four-manifold with one boundary component, identified with
\(S^3/\Gamma_n\). Suppose its metric has an outward Euclidean conical
collar: near the boundary it is \(dr^2+r^2g_{S^3/\Gamma_n}\), with \(r\)
increasing outward.

There are \textbf{no} conditions on the interior fundamental group,
intersection form, complex structure, spinor zeros, or the number of
smaller scales inside \(\Omega\). Then

\[
 \boxed{\int_\Omega\mathscr Q(g)\ge
              8\pi^2\frac{n-1}{2n}\ge\frac{8\pi^2}{3}.}
                                                        \tag{1.2}
\]

\hypertarget{boundary-formula-including-its-sign}{%
\subsubsection{1.1. Boundary formula, including its
sign}\label{boundary-formula-including-its-sign}}

Attach the exactly flat exterior cone to obtain a smooth complete ALE
manifold with compactly supported curvature. Its Euler characteristic
and signature are those of \(\Omega\). Use the convention

\[
 \chi=\frac1{8\pi^2}\int
       \left(|W^+|^2+|W^-|^2-\frac12|\operatorname{Ric}^{\circ}|^2
                                      +\frac{s^2}{24}\right)+\frac1n,
\] \[
 \tau=\frac1{12\pi^2}\int(|W^+|^2-|W^-|^2)
                                      +\eta_{\rm LV}(\Gamma_n).
                                                        \tag{1.3}
\]

Here \(\eta_{\rm LV}\) is the positive cotangent-sum convention of
\href{https://msp.org/gt/2016/20-3/gt-v20-n3-p13-s.pdf}{Lock--Viaclovsky,
equations (1-5), (3-1), and (3-5)}. It is the negative of the boundary
eta convention in which the signature formula has a minus sign. The
convention is fixed by (1.3), not by an unqualified occurrence of the
word eta.

For odd \(n\), the negative continued fraction is

\[
 \frac{n}{n-2}=[\underbrace{2,\ldots,2}_{(n-3)/2},3],
 \qquad (n-2)^{-1}\equiv\frac{n-1}{2}\pmod n.
\]

Substitution into the cyclic eta formula gives

\[
 \eta_{\rm LV}(\Gamma_n)
       =-\frac{(n-1)(n-5)}{6n},\qquad
 B_n:=\frac2n+3\eta_{\rm LV}(\Gamma_n)
       =3-\frac n2-\frac1{2n}.
                                                        \tag{1.4}
\]

Consequently the exact signed relation is

\[
 2\chi(\Omega)+3\tau(\Omega)-B_n
 =\frac1{4\pi^2}\int_\Omega
       \left(2|W^+|^2+\frac{s^2}{24}
                                  -\frac{|\operatorname{Ric}^{\circ}|^2}{2}\right).
                                                        \tag{1.5}
\]

This formula loses the whole \(W^-\) term. No upper bound for that term
has been inserted. Since the boundary is a rational homology sphere,
\(\tau(\Omega)\) is the integer signature of the relative intersection
pairing. Also

\[
 B_n=\frac{5-n}{2}+\frac{n-1}{2n},\qquad
 \operatorname{dist}(B_n,\mathbb Z)=\frac{n-1}{2n}\ge\frac13.
                                                        \tag{1.6}
\]

The absolute value of the integral in (1.5) is at most
\(\tfrac12\int\mathscr Q\). This proves (1.2).

For a sign check, \(n=3\) gives \(B_3=4/3\); \(n=5\) gives \(B_5=2/5\).
Reversing the cone orientation changes the calculation. In particular
the orientation reversal at \(n=3\) is an \(A_2\) cone and does have a
Ricci-flat hyperkähler filling. It is not covered by (1.2).

\hypertarget{actual-approximately-conical-collars}{%
\subsubsection{1.2. Actual approximately conical
collars}\label{actual-approximately-conical-collars}}

The assertion is stable under the collar hypothesis used in the
preceding annular theorem. More precisely, suppose an inner collar of
\(\Omega\) is identified with \(\{1/2<|y|\le1\}/\Gamma_n\), and the
\textbf{lift} of \(r^{-2}g\) has Euclidean \(C^2\) distance at most
\(\delta\). Then

\[
 \boxed{\int_\Omega\mathscr Q(g)
          \ge 8\pi^2\frac{n-1}{2n}-C\delta^2/n.}
                                                        \tag{1.7}
\]

The constant is independent of odd \(n\) and of the physical scale
\(r\). To prove this, leave the metric unchanged below radius \(3/4\),
and on the fixed annulus interpolate, with a fixed radial cutoff, to the
flat cone metric before reaching radius 1. The lifted interpolated
metric is positive definite for small universal \(\delta\), and its
curvature is \(O(\delta)\) in the rescaled coordinates. Its volume is
\(O(1/n)\). Thus the added integral of (1.1) is at most \(C\delta^2/n\).
Extend by the flat cone and apply (1.2). The unchanged interior
contributes its original integral; the nonnegative original collar
contribution can simply be retained in the upper bound. Four-dimensional
curvature-square integrals are invariant under the homothety, proving
(1.7) at every \(r\).

Fixing a sufficiently small universal \(\delta_0\), (1.7) in particular
gives the order-independent lower bound

\[
                  \int_\Omega\mathscr Q\ge q_0,
                  \qquad q_0=4\pi^2/3.              \tag{1.8}
\]

Unlike the finite-family constants in the spinor annular compactness
argument, this constant does not deteriorate as the quotient order
grows. The lifted metric control is essential; a bare diffeomorphism of
the intrinsic boundary to a lens space does not give (1.7).

\hypertarget{consequence-for-the-original-integral-defect}{%
\subsubsection{1.3. Consequence for the original integral
defect}\label{consequence-for-the-original-integral-defect}}

Let \(M\) be closed and oriented, let \(\omega\) be a smooth self-dual
two-form with \(|\omega|\le L\), and \(0\le t\le1/4\). As before set

\[
 A_{\omega,t}=t\left(\frac{|\omega|^2}{3}I-\omega\otimes\omega\right),
\] \[
 Z_t=|\operatorname{Ric}+g|^2+|W^+-A_{\omega,t}|^2
                         +( |\omega|^2-2)^2+|d\omega|^2,
 \qquad z=V^{-1}\int_M Z_t.                         \tag{1.9}
\]

Suppose \(\Omega_1,\ldots,\Omega_N\subset M\) are embedded smooth
one-boundary cores of the kind in Section 1.2, with overlap at most
\(\Lambda\). Their odd orders may vary without bound. If

\[
 \operatorname{Vol}(\Omega_i)\le v_0(L),\qquad
 v_0(L)=\frac{q_0}{2(8+L^4/3)},
                                                        \tag{1.10}
\]

then

\[
                    \boxed{\frac NV\le\frac{16\Lambda}{q_0}\,z.}
                                                        \tag{1.11}
\]

Indeed, \(|A_{\omega,t}|^2=(2/3)t^2|\omega|^4\le L^4/24\), so

\[
 \mathscr Q\le |\operatorname{Ric}|^2+4|W^+|^2
          \le8Z_t+8+L^4/3.                         \tag{1.12}
\]

Sum (1.8), use the overlap for the defect term, and absorb the sum of
core volumes using (1.10). This proves (1.11). No pointwise Ricci bound,
noncollapse, diameter, Sobolev constant, first-order expansion, or
control of \(W^-\) is used.

For the actual weighted forms of the previous proof, their uniform
amplitude bound and \(\int Z_{1/\beta}\le C\eta V\) give

\[
                         N/V\le C\Lambda\eta.       \tag{1.13}
\]

Thus a positive density of small-volume smoothly filled resonant cores
cannot occur in a zero-defect sequence, however small their individual
radii and however long the inner scale separation. The local resonant
annuli from the preceding construction remain valid: they did not fill
their quotient apex. Equations (1.7)--(1.13) concern the \textbf{whole}
smooth filling, not an isolated terminal bubble or an unmarked cap.

The original averaged hypotheses have not been proved to produce these
embedded cores, their overlap bound, or their small volumes. This result
therefore removes an actual smooth-filling alternative when present; it
does not itself prove the unrestricted integral zero-end theorem.

\hypertarget{intermediate-cores-and-cancellation-across-scales}{%
\subsubsection{1.4. Intermediate cores and cancellation across
scales}\label{intermediate-cores-and-cancellation-across-scales}}

For a smooth region \(W\) with one outgoing conical boundary of type
\(\Gamma_{n_0}\) and \(r\) incoming conical boundaries of types
\(\Gamma_{n_1},\ldots,\Gamma_{n_r}\), the same calculation gives

\[
 2\chi(W)+3\tau(W)-B_{n_0}+\sum_{i=1}^r B_{n_i}
   =\frac1{4\pi^2}\int_W
     \left(2|W^+|^2+\frac{s^2}{24}
                         -\frac{|\operatorname{Ric}^{\circ}|^2}{2}\right).
                                                        \tag{1.14}
\]

This is the boundary transgression formula with the incoming shape
operators and boundary orientations reversed. It can also be verified by
gluing arbitrary smooth caps and subtracting (1.5); signature and Euler
characteristic are additive across rational homology three-spheres.

In particular a one-child flat neck has no cost: identical boundary
terms cancel. It is wrong to apply (1.2) independently to every vertex
of a bubble tree. When the entire core is smoothly filled, all internal
terms in (1.14) cancel and (1.5) for its outer boundary survives. A
multiscale decomposition cannot remove the full-core cost.

There is also a uniform intermediate restriction. With two incoming
resonant boundaries, write \(f_n=(n-1)/(2n)\in[1/3,1/2)\). Then

\[
 f_{n_0}-f_{n_1}-f_{n_2}\in(-2/3,-1/6),\qquad
 \operatorname{dist}(B_{n_0}-B_{n_1}-B_{n_2},\mathbb Z)>1/6.
\]

Consequently an exactly conical binary branching region pays
\(\int_W\mathscr Q>4\pi^2/3\). This is not an arbitrary-valence rule:
with outgoing order 3 and three incoming orders 9 the fractional terms
differ by \(1/3-3(4/9)=-1\). The numerical obstruction then vanishes.
That arithmetic example asserts neither existence of a zero-cost
cobordism nor a solution of the actual field equations.

\hypertarget{the-full-core-cost-is-sharp-and-its-sign-has-content}{%
\subsubsection{1.5. The full-core cost is sharp, and its sign has
content}\label{the-full-core-cost-is-sharp-and-its-sign-has-content}}

The boundary value in (1.2) is the exact infimum, not only a convenient
lower bound:

\[
 \boxed{\inf_{\text{smooth conically collared fillings of }\Gamma_n}
                   \int\mathscr Q=4\pi^2\frac{n-1}{n}.}        \tag{1.15}
\]

The same infimum is obtained if the smooth topology is fixed to a
compact truncation of the minimal complex resolution. Indeed, its
exceptional chain is the negative definite chain \([2,\ldots,2,3]\) of
length \((n-1)/2\). Hence

\[
 \chi=(n+1)/2,\qquad \tau=-(n-1)/2,\qquad
                 2\chi+3\tau=(5-n)/2=\lfloor B_n\rfloor.
\]

\href{https://arxiv.org/pdf/math/0206229}{Calderbank--Singer, Theorem A}
provides a complete scalar-flat Kähler ALE metric on this resolution. In
the complex orientation \(W^+=s=0\), so (1.5), now taken on its complete
ALE end, gives

\[
 \int|\operatorname{Ric}^{\circ}|^2
                 =8\pi^2\frac{n-1}{2n}.             \tag{1.16}
\]

Truncate at radius \(R\to\infty\). The ALE estimates make the lifted
rescaled metric \(C^2\)-close to the flat cone on the truncation collar.
Flatten that collar by the construction of Section 1.2. The change in
the curvature-square integral tends to zero. Rescale the truncation to a
fixed boundary radius; the energy is unchanged. This gives the upper
bound in (1.15); (1.2) gives the lower bound. Further homotheties make
the core volume as small as desired without changing this cost. No
attainment assertion for the exactly conical compact problem is needed.

This uses a classical existence theorem to calibrate the new full-core
inequality; it is not a new scalar-flat ALE construction.

More precisely, put \(f_n=(n-1)/(2n)\) and
\(j=2\chi+3\tau-\lfloor B_n\rfloor\in\mathbb Z\). The signed identity is

\[
 P=4W_++S_2/12+8\pi^2(f_n-j),                       \tag{1.17}
\]

where the integrals are over the core. Thus \(j\le0\) forces
\(P\ge8\pi^2f_n\), whereas \(j\ge1\) forces
\(4W_++S_2/12\ge8\pi^2(1-f_n)\). The integer topology decides which side
pays. For a complete Ricci-flat ALE filling this implies
\(W_+\ge2\pi^2(1-f_n)=\pi^2(1+1/n)\); it does not rule out every non-ASD
Ricci-flat filling.

The sharp examples also delimit what this argument can control. On their
complete scalar-flat Kähler models,

\[
 \int|W^-|^2=4\pi^2n+6\pi^2-10\pi^2/n,
 \qquad \int\mathscr Q=4\pi^2(1-1/n).
                                                        \tag{1.18}
\]

Thus Euler number and opposite-chiral energy can grow without bound
while the full-core cost stays bounded. There can be no
order-independent estimate of core Euler number by this cost alone.
These are actual metrics on noncompact resolutions, not a closed
zero-end sequence with actual monopole supply. In particular (1.18) does
not contradict (1.11), which counts cores rather than their Euler
numbers.

\hypertarget{a-defect-atom-not-only-a-curvature-atom}{%
\subsubsection{1.6. A defect atom, not only a curvature
atom}\label{a-defect-atom-not-only-a-curvature-atom}}

The sign in (1.17) improves the translation to the actual defect. Write
\(e=\int_\Omega Z_t\), \(v=\operatorname{Vol}(\Omega)\), and
\(b_L=2/3+L^4/12\). The collar-flattening construction also controls the
\textbf{signed} error in (1.5): if its error in the unnormalized
integral is at most \(\epsilon=C\delta^2/n\), then

\[
 j\le0\quad\Longrightarrow\quad
 \int_\Omega|\operatorname{Ric}+g|^2
              \ge8\pi^2(f_n-j)-2\epsilon;                       \tag{1.19}
\] \[
 j\ge1\quad\Longrightarrow\quad
 e\ge\frac12\left[
       \sqrt{(4\pi^2(j-f_n)-\epsilon)_+}-\sqrt{b_Lv}
                         \right]_+^2.                          \tag{1.20}
\]

For (1.19), \(P\le\int|\operatorname{Ric}+g|^2\le e\) and (1.17)
suffice; no small-volume assumption is needed. For (1.20), apply the
Hilbert-space triangle inequality to \((\sqrt2W^+,s/\sqrt{24})\),
splitting it into the error
\((\sqrt2(W^+-A_{\omega,t}),(s+4)/\sqrt{24})\) and the bounded model
\((\sqrt2A_{\omega,t},-4/\sqrt{24})\). Since
\(|s+4|^2\le4|\operatorname{Ric}+g|^2\), this gives

\[
 \int_\Omega(2|W^+|^2+s^2/24)
                 \le(\sqrt{2e}+\sqrt{b_Lv})^2.
\]

Its left side is at least \(4\pi^2(j-f_n)-\epsilon\) by the signed
formula and \(P\ge0\). Thus (1.20) follows.

For any sequence of such smooth cores with \(v\to0\) and \(\delta\to0\),
possibly with unbounded odd orders,

\[
                       \boxed{\liminf\int_\Omega Z_t\ge\pi^2.} \tag{1.21}
\]

For fixed order the lower bound is \(\pi^2(1+1/n)\), or the stronger
\(4\pi^2(1-1/n)\) on the branch \(j\le0\). Unlike the sharp
\(\mathscr Q\) infimum, sharpness of (1.20)--(1.21) for an actual
monopole field is not asserted. An orbifold core with an unfilled apex
is not smooth and may have zero cost; its missing boundary contribution
must not be discarded.

There is an immediate pure-Ricci application: a bounded-overlap
collection of minimal-resolution cores (or any of the \(j\le0\)
topologies), with the prescribed small collar error, satisfies

\[
 \sum_i(f_{n_i}-j_i)\le C\Lambda
                         \int_M|\operatorname{Ric}+g|^2.        \tag{1.22}
\]

Absorb the collar errors into \(f_{n_i}-j_i\ge1/3\). This needs no
volume bound on the cores and no gauge field. Extra interior blow-ups
decrease \(j_i\) and incur additional linear cost, even if hidden at
smaller scales. The Calderbank--Singer truncations followed by shrinking
homotheties make \(\int|\operatorname{Ric}+g|^2\to8\pi^2f_n\): the
scalar cross term and the volume term vanish under that last homothety.
Thus the minimal-resolution pure-Ricci cost is also asymptotically sharp
as a metric-cap statement.

\hypertarget{a-second-obstruction-a-faithful-determinant-cannot-extend-flatly}{%
\subsection{2. A second obstruction: a faithful determinant cannot
extend
flatly}\label{a-second-obstruction-a-faithful-determinant-cannot-extend-flatly}}

The curvature cost has an independent global interpretation which avoids
the classification of Ricci-flat ALE spaces.

\textbf{Proposition.} Let \(X^m\), \(m\ge3\), be a smooth connected
complete ALE manifold with \(\operatorname{Ric}\ge0\) and nontrivial
infinity group \(\Gamma\). No flat unitary line bundle on \(X\)
restricts to a faithful holonomy character of \(\Gamma\).

\textbf{Proof.} Put \(N=|\Gamma|\); the asymptotic volume ratio is
\(1/N\). Choose a lift \(\tilde p\) of \(p\) in the universal cover and
a Dirichlet fundamental domain \(D\). The subset \(D\cap B(\tilde p,R)\)
projects, up to measure-zero cut loci, onto \(B_X(p,R)\). For a finite
subset \(S\) of deck transformations, let \(a\) be the largest
displacement of \(\tilde p\). The translates of this subset have
disjoint interiors and lie in \(B(\tilde p,R+a)\), so Bishop--Gromov
gives

\[
 |S|\operatorname{Vol}B_X(p,R)
     \le\operatorname{Vol}B_{\widetilde X}(\tilde p,R+a)
     \le\omega_m(R+a)^m.
\]

Letting \(R\) tend to infinity proves \(|S|\le N\). Thus the deck group
is finite of order \(d\le N\). A global flat character whose restriction
to \(\Gamma\) is faithful makes the inclusion-induced map
\(\Gamma\to\pi_1(X)\) injective. Hence \(d\ge N\) and \(d=N\).

For a finite cover the asymptotic volume ratio multiplies by its degree:
one may see this directly by taking the preimage of large balls, whose
distances from a fixed lift differ by a uniformly bounded amount because
there are finitely many deck transformations. Therefore
\(\operatorname{AVR}(\widetilde X)=d/N=1\). Equality in Bishop--Gromov
forces \(\widetilde X=\mathbb R^m\). A finite group of Euclidean
isometries fixes the average of any orbit and cannot act freely unless
it is trivial. This contradicts \(N>1\). QED.

The ALE covering facts can alternatively be obtained from
\href{https://arxiv.org/pdf/1102.2672}{Șuvaina, Lemmas 4.1--4.2 and
Remark 4.3}. The argument above displays the volume reason, not merely
an appeal to the Kähler ALE classification.

\hypertarget{application-to-the-actual-positive-spinor-alone}{%
\subsubsection{2.1. Application to the actual positive spinor
alone}\label{application-to-the-actual-positive-spinor-alone}}

There is no smooth complete Ricci-flat ALE four-manifold with oriented
infinity \(1/n(1,-2)\), odd \(n>1\), carrying a nonzero global parallel
positive Spin-c spinor. The opposite twisted spinor need not extend.

First, flatness of the determinant is automatic. For a parallel spinor
\(\Phi\) with determinant curvature \(F_A=i\Omega\), the local curvature
identity is

\[
 \operatorname{Ric}(v)\cdot\Phi
                      =i(v\mathbin{\lrcorner}\Omega)\cdot\Phi.
                                                        \tag{2.1}
\]

See \href{https://moroianu.perso.math.cnrs.fr/tex/1997cmp.pdf}{Moroianu,
Lemma 3.1}. For Ricci-flat metrics the right side vanishes. Clifford
multiplication by a nonzero real covector is invertible, and
\(\Phi\ne0\); hence \(\Omega=0\).

At infinity the parallel spinor gives an invariant vector of the
limiting positive representation. In the stated marking the canonical
positive weights are \((1,\zeta^{-1})\). A character twist permits an
invariant vector only for twists \(1\) and \(\zeta\); its determinant
character is respectively \(\zeta^{-1}\) and \(\zeta\). Both are
faithful. The proposition applies. A Ricci-flat Kähler ALE metric in
this orientation has such a parallel canonical Spin-c spinor, so it is
excluded as well. For trivial infinity group, Bishop--Gromov already
gives flat Euclidean space.

This deduction removes the first-order expansion, first-order curvature
energy, same-scale derivative bound, and global parallel-frame premises
from this particular resonant-filling exclusion. It still requires a
\textbf{genuine smooth complete ALE limit} and a positive spinor that
extends parallel over that same limit. Neither follows just by
inspecting an exterior annulus.

An independent consequence of Section 1 is the nonexistence of any
smooth Ricci-flat anti-self-dual ALE filling of this oriented cone, even
without spinor data. This agrees with, but does not require, the
Ricci-flat Kähler classification: its nontrivial cyclic types have
\(N=dk^2\), \(q=dkm-1\), whose determinant character has order \(k\),
not the order \(N\) demanded by \(q=-2\).

\hypertarget{an-integral-lower-ricci-bridge-in-an-elliptic-class}{%
\subsection{3. An integral lower-Ricci bridge in an elliptic
class}\label{an-integral-lower-ricci-bridge-in-an-elliptic-class}}

There is a second way to extend the endpoint theorem. Here the
additional information is an elliptic equation for Ricci, not a marked
core. The statement is not a claim that arbitrary small integral Ricci
error has the following regularity.

Fix \(p>2\), \(r_0\in(0,1]\), \(A,L<\infty\), and \(\Lambda_0\). Define
the negative part of the least Ricci eigenvalue by

\[
 \rho_-(x)=\max\{0,-\lambda_{\min}(\operatorname{Ric}^{\sharp}(x))\},
 \qquad k_p(r)=r^2\sup_x
                  \left(\fint_{B_r(x)}\rho_-^p\right)^{1/p}.
                                                        \tag{3.1}
\]

There is \(\delta_p>0\) such that the hypotheses

\[
 |\nabla^*\nabla\operatorname{Ric}|
              \le A|\operatorname{Rm}|\,|\operatorname{Ric}|,
 \quad k_p(r_0)\le\delta_p,
 \quad |\omega|\le L,
 \quad 8\pi^2\chi/V\le\Lambda_0                       \tag{3.2}
\]

on a closed connected smooth oriented four-manifold imply

\[
                  \boxed{\fint_M Z_t\ge c(p,r_0,A,L,\Lambda_0)>0,}
                  \qquad 0\le t\le1/4.             \tag{3.3}
\]

In this section \(|\operatorname{Rm}|\) is the full four-index tensor
norm. No diameter, pointwise Ricci, or noncollapse bound is assumed.
Both the elliptic condition and the \textbf{small} normalized local
lower-Ricci budget in (3.2) are genuine inputs.

\hypertarget{the-normalized-local-estimate}{%
\subsubsection{3.1. The normalized local
estimate}\label{the-normalized-local-estimate}}

\href{https://web.math.ucsb.edu/~wei/paper/Dai-Wei-Zhang.pdf}{Dai--Wei--Zhang,
Theorem 1.1, Corollary 1.5, and (2.5)--(2.6)} give, after choosing
\(\delta_p\) small enough and rescaling,

\[
 \frac{\operatorname{Vol}B_a(x)}{\operatorname{Vol}B_b(x)}
                 \ge\frac12(a/b)^4,
                 \qquad 0<a\le b\le r_0.           \tag{3.4}
\]

The propagation to every smaller scale uses their (2.6); decrease the
initial threshold by the factor needed there. In particular

\[
 \operatorname{Vol}B_r\le2\omega_4r^4,
 \qquad \operatorname{Vol}B_{2r}\le32\operatorname{Vol}B_r.
                                                        \tag{3.5}
\]

When \(2r\le r_0\) and \(B_{2r}(x)\) is proper, apply their Dirichlet
Sobolev estimate at reference scale \(2r\), to \(|f|^3\), and then use
Cauchy--Schwarz. It yields

\[
 \left(\fint_{B_r}|f|^4\right)^{1/2}
            \le S(p)r^2\fint_{B_r}|\nabla f|^2,
                         \qquad f\in C_c^\infty(B_r).           \tag{3.6}
\]

This is a normalized Sobolev estimate. No uniformly bounded unnormalized
Sobolev constant has been introduced.

Under (3.2)'s elliptic inequality, (3.6) implies the following local
lemma, for constants \(e_0,C_R>0\) depending only on \(p,A\):

\[
 r^4\fint_{B_r}|\operatorname{Rm}|^2\le e_0
 \quad\Longrightarrow\quad
 \sup_{B_{r/8}}|\operatorname{Ric}|
                  \le C_R\left(\fint_{B_r}|\operatorname{Ric}|^2\right)^{1/2}.
                                                        \tag{3.7}
\]

Here are the normalization and proof, rather than an invocation of a
collapsed harmonic-radius theorem. Work with \(h=r^{-2}g\) and the
probability measure \(d\nu=dV_g/\operatorname{Vol}_g(B_r)\). A constant
multiple of Riemannian measure preserves integration by parts. With
\(U=\operatorname{Ric}_h\), \(B=\operatorname{Rm}_h\), all norms below
use this same measure on nested subballs of \(B_1^h\). The rescaled
Sobolev inequality is \(\|f\|_4^2\le S\|\nabla_hf\|_2^2\).

Kato's inequality, the tensor energy identity, and a cutoff yield

\[
 \|U\|_{4;B_{3/4}}^2+\|\nabla U\|_{2;B_{3/4}}^2
                              \le C\|U\|_{2;B_1}^2,             \tag{3.8}
\]

after absorbing \(C\|B\|_2\|\phi U\|_4^2\). Next use the universal
curvature identity \(\Delta B=\mathcal L(\nabla^2U)+B*B\). One
integration by parts and Young's inequality give

\[
 \int\psi^2|\nabla B|^2\le
 C\int(|\nabla\psi|^2|B|^2+\psi^2|\nabla U|^2)
                                      +C\int\psi^2|B|^3.
\]

Sobolev and (3.8) absorb the cubic term, because
\(\int\psi^2|B|^3\le\|B\|_2\|\psi B\|_4^2\). Thus

\[
                         \|B\|_{4;B_{1/2}}^2\le C\|B\|_{2;B_1}^2.
                                                        \tag{3.9}
\]

Finally \(u=|U|\) obeys \(\Delta u\ge-A|B|u\), where here
\(\Delta=\operatorname{tr}\nabla^2\) has nonpositive spectrum. For the
Moser test function \(v=\psi u^{m/2}\),

\[
 \int|B|v^2\le\|B\|_4\|v\|_{8/3}^2
                   \le\|B\|_4\|v\|_2\|v\|_4.
\]

Young's inequality and (3.6) give, on nested balls,
\(\|u\|_{2m,B_a}\le[C\{m^4+m^2(b-a)^{-2}\}]^{1/m} \|u\|_{m,B_b}\).
Iteration proves (3.7). These are the early integration-by-parts
estimates of \href{https://arxiv.org/pdf/math/0310302}{Tian--Viaclovsky,
Section 3, Lemmas 3.2--3.5}, written with the normalized measure; no
later volume-comparison premise from that paper is being transplanted.

\hypertarget{euler-density-selects-enough-good-centers}{%
\subsubsection{3.2. Euler density selects enough good
centers}\label{euler-density-selects-enough-good-centers}}

Assume \(z\le1\), since otherwise (3.3) is immediate. Then

\[
 \fint|\operatorname{Ric}|^2\le10,\qquad
 \fint|\operatorname{Rm}|^2
       =32\pi^2\chi/V+4\fint|\operatorname{Ric}^{\circ}|^2
       \le H:=4(1+\max\{\Lambda_0,0\}).             \tag{3.10}
\]

For any nonnegative \(F\), doubling and Fubini imply

\[
                     \fint_M\fint_{B_r(x)}F\le32\fint_M F.      \tag{3.11}
\]

Indeed, if \(x\in B_r(y)\), then
\(\operatorname{Vol}B_r(y)\le\operatorname{Vol}B_{2r}(x) \le32\operatorname{Vol}B_r(x)\).

Decrease \(e_0\le1\) so that \(2\omega_4e_0\) is below the Cheeger--Tian
local covering energy threshold for dimensionless Ricci bound \(C_R+1\),
allowing for the fixed norm conversion. Choose

\[
 R=\min\{r_0/64,1,(e_0/(2^{20}H))^{1/4}\}.           \tag{3.12}
\]

If \(D=\operatorname{diam}M\ge8R\), all \(B_{2R}(x)\) are proper. At
least half the centers satisfy
\(R^4\fint_{B_R(x)}|\operatorname{Rm}|^2\le e_0\), by (3.10)--(3.11). On
each such center (3.7) gives a pointwise Ricci bound
\(C_R\sqrt{e_0}R^{-2}\) on \(B_{R/8}\), while (3.5) gives
\(\int_{B_R}|\operatorname{Rm}|^2\le2\omega_4e_0\).

Apply the original bounded-Ricci local covering result of Cheeger--Tian
on the buffered ball \(B_{R/16}\), exactly as verified in
\href{20260927_local_geometry_and_variational_extrema.md}{I45, Section
3}. This gives fixed-size uniformly regular covering charts over a
positive volume fraction of centers. I45's local coercivity for \(Z_t\)
and its Liouville-measure averaging now give (3.3). Radial Jacobian
comparison is used only inside these good neighborhoods, where Ricci has
just been bounded; extension of a nonnegative angular integral to all
centers does not require a global Ricci bound.

If \(D<8R\), set \(u=D/16\). Every \(B_{2u}(x)\) is proper and
\(B_{2D}(x)=M\). Since \(2D<r_0\), (3.4) gives
\(\operatorname{Vol}B_u(x)\ge V/(2\cdot32^4)\). Therefore every center
satisfies

\[
 u^4\fint_{B_u(x)}|\operatorname{Rm}|^2
              \le32D^4H\le2^{17}R^4H\le e_0/8.
\]

By (3.7) and (3.10),
\(|\operatorname{Ric}|(x)\le C_R\sqrt{20\cdot32^4}\) everywhere. I45 now
applies directly, with a bound independent of the small diameter. This
completes the proof of (3.3).

\hypertarget{actual-and-variational-interfaces}{%
\subsubsection{3.3. Actual and variational
interfaces}\label{actual-and-variational-interfaces}}

The actual form supply gives \(\int Z_{1/\beta}\le C\eta V\), hence
(3.3) yields the uniform actual endpoint gap for all \(\beta\ge4\) under
(3.2). The earlier small-defect identity that makes \(V\) comparable to
\(T_+\) is purely algebraic, so a bounded \(\chi/T_+\) may replace the
Euler-density premise here as well.

For a constant-scalar-curvature critical point of
\(Q_\tau=\int|\operatorname{Ric}|^2+\tau\int s^2\), the equation is

\[
 \Delta\operatorname{Ric}
 =-2\operatorname{Rm}(\operatorname{Ric})
   -2\tau s\operatorname{Ric}
                   +\tfrac12(|\operatorname{Ric}|^2+\tau s^2)g.
                                                        \tag{3.13}
\]

Thus (3.2)'s elliptic inequality holds with \(A\le C(1+|\tau|)\). See
\href{https://arxiv.org/abs/1404.0569}{Catino, Proposition 2.1}. The
sharp-functional parameter \(\tau=-1/(4-\alpha)\) for \(0\le\alpha<1\)
lies in a fixed bounded interval. This establishes an actual connection
with the extremal problem, not an extra PDE silently imposed on
arbitrary near-Einstein metrics. A Yamabe representative only has
constant scalar curvature; it need not satisfy (3.13), and a minimizing
sequence need not be a critical sequence.

The small local quantity \(k_p(r_0)\) is not supplied by the global
\(L^2\) defect. Indeed, the comparison theorem itself requires
smallness, not merely a finite bound; Dai--Wei--Zhang's Remark 1.2
explicitly records this distinction. For an elliptic actual endpoint
sequence escaping (3.3), for every fixed \(p>2\) and \(0<r\le1\),

\[
                         \liminf_j k_p^{g_j}(r)\ge\delta_p.      \tag{3.14}
\]

A diagonal subsequence therefore has radii \(r_j\downarrow0\) and
centers \(x_j\) with
\(r_j^{2p}\fint_{B_{r_j}(x_j)}\rho_{-,j}^p \ge(\delta_p/2)^p\). This is
normalized local Ricci concentration, not an unnormalized energy atom or
an automatically existing ALE bubble.

\hypertarget{actual-spinor-remainders-and-concentration-of-smooth-extrema}{%
\subsection{4. Actual spinor remainders and concentration of smooth
extrema}\label{actual-spinor-remainders-and-concentration-of-smooth-extrema}}

This section concerns the remaining \textbf{coefficient}, not the
already settled linear order. All metrics here have constant scalar
curvature \(s=-4\). Ordinary SW solvability is required for the one
solution used below. It follows from the project's conformal-weight
supply by taking weight 1; no comparison of two different solutions is
made.

\hypertarget{an-exact-nonnegative-decomposition-at-the-ordinary-solution}{%
\subsubsection{4.1. An exact nonnegative decomposition at the ordinary
solution}\label{an-exact-nonnegative-decomposition-at-the-ordinary-solution}}

Let \(D_{\mathcal A}\Phi=0\), \(-iF_{\mathcal A}^+=\sigma(\Phi)\), with
\(|\sigma(\Phi)|^2=|\Phi|^4/8\). Put \(\rho=|\Phi|^2\), \(x=\rho/4\),
\(\omega=\sigma(\Phi)\), and
\(\gamma=\operatorname{Im}\langle\Phi,\nabla_{\mathcal A}\Phi\rangle\).
The maximum principle gives \(0\le x\le1\). Define

\[
 \begin{aligned}
 K&=\int|\nabla_{\mathcal A}\Phi|^2,&Q&=\int(4-\rho)^2,\\
 \Gamma&=\int(|d\rho|^2/4+|\gamma|^2),&
 J&=\|F_{\mathcal A}^+-(F_{\mathcal A}^+)_H\|_2^2,\\
 W_0&=\tfrac12(|\omega|^2I/3-\omega\otimes\omega),&
 F&=\int|W^+-W_0|^2,\\
 H&=\tfrac23\int(1-x)^2(2+4x-x^2).&&
 \end{aligned}                                                   \tag{4.1}
\]

The harmonic projection uses the same curvature/Chern sign as the
solution; its squared norm is \(4\pi^2p\). For
\(D=D_2=24V+18\int|W^+|^2-72\pi^2p\), one has the exact identity

\[
                  \boxed{D/18=F+\Gamma/2+J+H.}                  \tag{4.2}
\]

In particular \(F,J\le D/18\), \(\Gamma\le D/9\), and \(Q\le2D/3\). This
is a retention of the remainders at one actual ordinary solution, not a
new nonvanishing theorem. The underlying equations and Weitzenböck
conventions are those of
\href{https://arxiv.org/pdf/math/0003068}{LeBrun, equations (5)--(7)}.

For completeness, the Hopf differential and Dirac equation give

\[
 |\nabla\omega|^2=\frac\rho2|\nabla_{\mathcal A}\Phi|^2
                                      -\frac12|\gamma|^2,
 \qquad |d\omega|^2=|\gamma|^2/4.                  \tag{4.3}
\]

These factors can be checked in the Clifford frame \(C_1=I\),
\(C_2=\operatorname{diag}(i,-i)\), \(C_3=((0,1),(-1,0))\),
\(C_4=((0,-i),(-i,0))\). At \(\Phi=\sqrt\rho\,e_1\), write
\(\nabla_i\Phi=(u_i+iv_i,w_i+iz_i)\). Differentiating in the self-dual
basis \((e^{12}+e^{34},e^{13}-e^{24},e^{14}+e^{23})\) gives
\(\nabla_i\omega=(\sqrt\rho/2)(u_i\Sigma_1+z_i\Sigma_2-w_i\Sigma_3)\).
The two Dirac equations give \(d^*\omega=-\gamma/2\), with an irrelevant
simultaneous sign change if the opposite curvature convention is used.
The squared statements extend across zeros polynomially.

Pair the scalar spinor equation with \(\rho\) and integrate to obtain

\[
 2\int|d\rho|^2+4\int\rho|\nabla_{\mathcal A}\Phi|^2
                              =\int(4-\rho)\rho^2.
\]

Thus, writing \(T_\omega=\int(|\nabla\omega|^2-2|d\omega|^2)\),

\[
 T_\omega=8\int x^2(1-x)-\Gamma,
 \qquad
 F=\int|W^+|^2+\tfrac23\int x^4+T_\omega/2-\tfrac43\int x^2.
                                                        \tag{4.4}
\]

Also \(J=2\int x^2-4\pi^2p\). Substitution reduces (4.2) to
\((1-x)^2(2+4x-x^2)=2-7x^2+6x^3-x^4\). Finally \(Q/12\le H\le5Q/24\),
proving the stated inequalities.

\hypertarget{the-exact-critical-equation-and-cotton-identity}{%
\subsubsection{4.2. The exact critical equation and Cotton
identity}\label{the-exact-critical-equation-and-cotton-identity}}

Write \(Z=\operatorname{Ric}^{\circ}\), \(P=\int|Z|^2>0\), and
\(\alpha=\mathcal E(g)=4P/(P+4V)\). Set

\[
 \tau_*=-1/(4-\alpha),\qquad
 b=\alpha/(4-\alpha)=P/(4V),\qquad q=1/3-2b.
\]

Here \(\Delta=\operatorname{tr}\nabla^2\) has nonpositive spectrum; this
is the opposite sign to the positive Laplacian used for the earlier
conformal variation. Use \(R_{ikjl}g^{kl}=\operatorname{Ric}_{ij}\),
\(\sec(e_i,e_j)=R_{ijij}\). The genuine \(L^2\) metric gradient of
\(Q_{\tau_*}\), holding its parameter fixed while varying the metric, is

\[
 \mathcal G=-\Delta Z-2\operatorname{Rm}(Z)-2bZ
                                      +\tfrac12|Z|^2g.          \tag{4.5}
\]

This follows from the classical first variation (3.13). In particular, a
smooth attained positive minimum of \(\mathcal E\) has \(\mathcal G=0\)
by the preceding variational reduction. Formula (4.5) itself is valid
without stationarity.

Define the Ricci-normalized Cotton tensor by
\(C_{ijk}=\nabla_kZ_{ij}-\nabla_jZ_{ik}\), with its full ordered tensor
norm; it is not divided by a Schouten factor. Integration by parts,
using \(\operatorname{div}Z=0\), gives

\[
 \tfrac12\int|C|^2
   =\int\langle\mathcal G,Z\rangle
     +\int\sum_{i<j}(\lambda_i+\lambda_j)^2\sec_{ij}+2bP,       \tag{4.6}
\]

where \(\lambda_i\) are the eigenvalues of \(Z\). Indeed the usual
Cotton identity is
\(\int|\nabla Z|^2-\tfrac12\int|C|^2 =\int(\operatorname{Rm}(Z,Z)-\operatorname{tr}Z^3-s|Z|^2/4)\);
pair (4.5) with \(Z\) and use

\[
 \operatorname{Rm}(Z,Z)+\operatorname{tr}Z^3+s|Z|^2/4
                  =\sum_{i<j}(\lambda_i+\lambda_j)^2\sec_{ij}.
\]

The source formulas are \href{https://arxiv.org/pdf/1404.0569}{Catino,
Propositions 2.1 and 4.1}, not his pinching theorem for a different
parameter interval.

Define \(K^\perp(\Pi)=(\sec\Pi+\sec\Pi^\perp)/2\) and
\(k_\perp=\max_\Pi K^\perp(\Pi)\). The complementary weights in (4.6)
are equal and their total is \(2|Z|^2\). For a stationary metric, for
every \(k>b\),

\[
 \int(k_\perp+k)_+|Z|^2
                  \ge(k-b)P+\tfrac14\int|C|^2.                 \tag{4.7}
\]

In particular a positive smooth extremizer cannot have \(k_\perp<-b\)
everywhere. No strict conclusion at the equality threshold is inferred
without additional information about Cotton.

\hypertarget{joint-use-of-the-actual-field-and-the-full-variational-equation}{%
\subsubsection{4.3. Joint use of the actual field and the full
variational
equation}\label{joint-use-of-the-actual-field-and-the-full-variational-equation}}

Let \(w_3^\pm\) be the largest Weyl eigenvalues. Decomposition into a
unit self-dual and anti-self-dual form, followed by complementary-plane
averaging, cancels the Ricci off-diagonal block and gives

\[
 k_\perp=-1/3+(w_3^++w_3^-)/2.
\]

Set \(\xi=(w_3^++w_3^--1/3)_+\). Equations (4.6)--(4.7) imply

\[
 qP+\tfrac12\int|C|^2
             \le\int\langle\mathcal G,Z\rangle+\int\xi|Z|^2.
                                                        \tag{4.8}
\]

The top eigenvalue of the actual \(W_0\) is \(x^2/3\le1/3\). For a
traceless symmetric three-by-three matrix,
\(\lambda_{\max}\le\sqrt{2/3}|\cdot|\). Therefore

\[
 \int\xi^2\le\tfrac89(F+3B)
            \le\tfrac89(D/18+3B),\qquad B=\int|W^-|^2,          \tag{4.9}
\]

using \((a+b)^2\le(4/3)(a^2+3b^2)\). The weighting matters here. Assume
the actual budget \(p\ge A>0\), \(T_+=A(1-h)\), and write
\(u_-=3T_-/A\), \(z_a=p/A-1\), \(e=R/(8\pi^2A)\). Chern--Weil gives the
exact identities

\[
 D/18=2\pi^2A(e-1-h-2z_a),\qquad
 B=P/3+(2\pi^2A/3)(u_--e).
\]

Consequently

\[
 \boxed{\mathcal H:=D/18+3B
             =P+2\pi^2A(u_--1-h-2z_a).}             \tag{4.10}
\]

Only when \(T_-\le T_+\) is \(u_-\) the earlier two-sided \(U\). This
holds eventually on the balanced small-h family; it is not assumed for
arbitrary orientations.

For a stationary metric with \(0<\alpha<4/7\), \(q>0\) and (4.8) forces
\(\mathcal H>0\). Hölder and (4.9) prove

\[
 \boxed{\int|Z|^4\ge\frac98q^2\frac{P^2}{\mathcal H},
 \qquad \|Z\|_\infty^2\ge\frac98q^2\frac P{\mathcal H}.}       \tag{4.11}
\]

More generally, without assuming stationarity, one has

\[
 \boxed{
 \frac{\|\mathcal G\|_2}{\sqrt P}
       +\sqrt{\frac{8\mathcal H}{9P}}\,\|Z\|_\infty
               \ge q+\frac{\int|C|^2}{2P}.}         \tag{4.12}
\]

One may replace \(\|Z\|_\infty\) by \((\int|Z|^4/P)^{1/2}\). This is
obtained directly from (4.8) by Cauchy--Schwarz; it remains valid when
\(\mathcal H=0\).

\hypertarget{explicit-consequence-on-the-realized-balanced-family}{%
\subsubsection{4.4. Explicit consequence on the realized balanced
family}\label{explicit-consequence-on-the-realized-balanced-family}}

Use the same actual family as I38--I41, with \(h_j\to0\) and
\(U_j=F(h_j)+o(h_j^2)=1+h_j+O(h_j^2)\). For any scalar \(-4\) competitor
satisfying \(0<\alpha_j=\mathcal E(g_j)\le4h_j/(1+h_j)\), the existing
actual linear theorem and \(e_j\ge1+h_j\) give \(P_j\ge cA_jh_j\). Also
\(b_j\le h_j\), and the favorable sign of \(z_a\ge0\) in (4.10) gives
\(\mathcal H_j/P_j\le1+O(h_j)\).

Thus if these competitors are stationary, in particular if they are
smooth attained minima,

\[
 \boxed{
 \liminf_j\frac{\int|\operatorname{Ric}^{\circ}_{g_j}|^4}{P_j}\ge\frac18,
 \qquad
 \liminf_j\|\operatorname{Ric}^{\circ}_{g_j}\|_{\infty,\rm HS}^2\ge\frac18.}
                                                        \tag{4.13}
\]

Since \(s=-4\), \(\operatorname{Ric}+g=Z\), and four-dimensional
Hilbert--Schmidt norm is at most twice the spectral norm. Consequently

\[
                         \liminf_j\delta_-(g_j)\ge\frac1{4\sqrt2}.
                                                        \tag{4.14}
\]

If instead these scalar \(-4\) competitors have \(\delta_-(g_j)\to0\),
then (4.12) forces the explicit variational residual

\[
                   \boxed{\liminf_j
                       \|\mathcal G_j\|_2/\sqrt{P_j}\ge1/3.}   \tag{4.15}
\]

They cannot be relatively almost critical for this quadratic-curvature
problem. This is an unrestricted metric variation, not a variation
confined to constant-scalar metrics. With
\(R=\int|\operatorname{Ric}|^2\) and \(S_2=\int s^2\), the
scalar-variance term has zero first variation at a constant-scalar
metric. Hence
\(\nabla\mathcal E=\nabla\mathcal J=(S_2/R^2)\mathcal G =((4-\alpha)/R)\mathcal G\).
In computing \(\mathcal G\), the parameter \(\tau\) is selected at the
metric and then held fixed; differentiating \(\tau(g)\) inside
\(Q_{\tau(g)}(g)\) would instead differentiate an identically zero
expression.

No existence of either sequence is asserted. These conclusions do not
prove nonattainment, determine the all-metric coefficient, or apply to
general nonconstant-scalar input after an uncontrolled conformal change.
They identify an actual obstruction to treating weak integral
minimization as a uniformly pointwise near-Einstein deformation problem.

\hypertarget{the-signed-term-still-needed-for-coefficient-four}{%
\subsubsection{4.5. The signed term still needed for coefficient
four}\label{the-signed-term-still-needed-for-coefficient-four}}

The ordinary solution also makes the unclosed estimate explicit. Define
\(L_0=2H-Q/8=(2/3)\int(1-x)^2(1+8x-2x^2)\ge0\). Retaining
\(\Sigma=8K+Q+8J\) and using (4.2) gives

\[
 \boxed{P-8\pi^2Ah
        =2[2F+\Gamma+J+L_0+4\pi^2(p-A)-K].}          \tag{4.16}
\]

The same actual solution supplies every term. Its only negative payment
on the right is the spinor gradient energy \(K\); nonnegative losses
must not be discarded before asking whether they compensate it. On
\(V-2\pi^2A=O(Ah)\), the estimate
\([K-\Gamma-J-L_0-4\pi^2(p-A)]_+=o(Ah)\) would suffice for leading
coefficient four, even if the earlier unsigned directional cost is only
\(O(Ah)\). This estimate is \textbf{not} proved here. Equations
(4.13)--(4.15) give information about extremal geometry instead of
pretending that the remaining signed balance has already closed.

\hypertarget{what-this-closes-and-what-still-requires-mathematics}{%
\subsection{5. What this closes, and what still requires
mathematics}\label{what-this-closes-and-what-still-requires-mathematics}}

The preceding local annulus theorem has two actual exceptional lifts.
Neither lift has been discarded. The new distinction is that both solve
the local field equations while their \textbf{entire smooth filling}
must pay the curvature cost (1.7). This is stronger than ruling out one
proposed first-order correction, and applies to arbitrary topology
inside a core with the specified outer collar.

It does not yet supply a decomposition of every weak-error sequence into
such cores. Hyperkähler-type, nonresonant annuli are still governed by
the actual-field estimates, not by (1.2); reversing orientations is
prohibited. No conclusion about all Einstein metrics on the fixed
Dehn-twist manifold is inferred here. The balanced family's sharp
coefficient and the pure integral-Ricci endpoint remain separate
problems.

The elliptic bridge in Section 3 replaces pointwise Ricci control by a
specific curvature equation and a small local normalized lower-Ricci
budget. It does not hide these additional hypotheses inside an
admissibility definition. Section 4 keeps the ordinary actual solution's
full remainders, then links them to the variational equation: small
total defect does not make a smooth low-energy extremum uniformly
pointwise near Einstein.

The next sufficient statements are now precise: obtain the marked
complete core structure from weaker original inputs, or estimate the
signed ordinary spinor balance in (4.16). Neither requires redoing
linear sharpness or repeating the local holonomy classification. The
proof above does not assume either missing statement.

END COMPLETE SOURCE

\begin{center}\rule{0.5\linewidth}{0.5pt}\end{center}

\hypertarget{source-2-20260927_signed_balance_and_conformal_flow.md}{%
\subsection{Source 2:
20260927\_signed\_balance\_and\_conformal\_flow.md}\label{source-2-20260927_signed_balance_and_conformal_flow.md}}

Original path:
continuation/20260927\_signed\_balance\_and\_conformal\_flow.md

SHA-256:
9d726a0c5653f223eedb424d985704918dd1fdf8ac0d01317d5d063a6cc39dd7

BEGIN COMPLETE SOURCE

\hypertarget{signed-spinor-energy-conformal-ricci-flow-and-extremal-degeneration}{%
\section{Signed spinor energy, conformal Ricci flow, and extremal
degeneration}\label{signed-spinor-energy-conformal-ricci-flow-and-extremal-degeneration}}

27 September 2026. Continuation from \texttt{754bc3a}, on the local
branch \texttt{local/yau26-signed-balance-20260927}. The preceding proof
is \href{20260927_full_core_cost_and_extremal_metrics.md}{Full-core cost
and extremal metrics}. This is a mathematical continuation, not a
web-export packet. The inherited linear small-h theorem is used, not
rediscovered. All statements below retain their own hypotheses; no
general zero-end gap or balanced coefficient-four theorem is asserted.

The main connection is now dynamic: the all-metric gap and the two
energy walls bound the lifetime of an actual geometric flow, even when
the gap itself tends to zero. A genuine ordinary-SW deformation also
determines which term cannot be removed from the signed sharpness
argument. The last sections identify the additional information supplied
by stationarity and by noncollapsed curvature blow-up.

\hypertarget{three-static-budgets-force-finite-conformal-flow-time}{%
\subsection{1. Three static budgets force finite conformal-flow
time}\label{three-static-budgets-force-finite-conformal-flow-time}}

Let \(M\) be closed, connected and oriented. Write

\[
 V=\operatorname{Vol}(g),\quad P=\int|\operatorname{Ric}^{\circ}|^2,
 \quad R=\int|\operatorname{Ric}|^2,
 \quad S_2=\int s^2.
\]

Suppose constants \(A>0\), \(h>0\), and \(0<\mu<1\) give, for every
smooth metric on \(M\),

\[
 S_2\ge32\pi^2A,\qquad R\ge8\pi^2A(1+h),
 \qquad \mathcal E(g)\ge\mu.                         \tag{1.1}
\]

Only the corresponding bounds on scalar-minus-four metrics are needed.
In the realized balanced surgery family, the actual gauge input supplies
the first two inequalities and I40 supplies \(\mu=c_1h\), after
decreasing a fixed \(c_1\) if necessary. Put

\[
 U=\frac{3(2\chi-3\tau)}A.
\]

This uses the specified opposite orientation, not an unlabelled minimum
over orientations. It agrees with the earlier \(U\) on the balanced
family for sufficiently small \(h\).

\hypertarget{the-actual-evolution-and-the-basic-lifetime-bound}{%
\subsubsection{1.1. The actual evolution and the basic lifetime
bound}\label{the-actual-evolution-and-the-basic-lifetime-bound}}

Start at any smooth metric with \(s=-4\). Use positive scalar Laplacian
\(\Delta=d^*d\) and conformal pressure \(\varpi\). Conformal Ricci flow
is

\[
 \partial_tg=-2Z-2\varpi g,\qquad
 Z=\operatorname{Ric}+g=\operatorname{Ric}^{\circ},\qquad
 (3\Delta+4)\varpi=|Z|^2.                            \tag{1.2}
\]

This is the classical scalar-preserving flow, not ordinary normalized
Ricci flow. Short-time smooth existence follows from
\href{https://arxiv.org/pdf/1109.5377}{Lu--Qing--Zheng, Theorem 1.1}:
the required elliptic operator is invertible at negative scalar
curvature. Their nonpositive-Laplacian convention is converted in (1.2).
The scalar variation verifies preservation of \(s=-4\), and integration
of the pressure equation gives

\[
 \varpi\ge0,\qquad 4\int\varpi=P,
 \qquad \boxed{V'=-P.}                              \tag{1.3}
\]

The volume identity is also equation (2.3) of that paper. It is a
classical input; the following use of the project budgets is the new
connection here.

At scalar \(-4\),

\[
 R=P+4V,\qquad \mathcal E=\frac{4P}{P+4V},\qquad
 V\ge V_*:=2\pi^2A.
\]

Consequently every smooth interval of (1.2) satisfies

\[
 V'\le-\frac{4\mu}{4-\mu}V,
 \qquad
 \boxed{T_{\max}\le\frac{4-\mu}{4\mu}
                      \log\frac{V(0)}{V_*}.}        \tag{1.4}
\]

Indeed integrate the differential inequality and then use the fixed
positive lower volume bound. Here \(T_{\max}\) is the maximal classical
smooth continuation time in normalization (1.2). The conclusion neither
identifies the singularity nor assumes a curvature-only extension
criterion. In particular it does not prove nonexistence by assuming a
flow must converge: its input is an already established all-metric gap.

\hypertarget{both-energy-walls-improve-the-time-bound}{%
\subsubsection{1.2. Both energy walls improve the time
bound}\label{both-energy-walls-improve-the-time-bound}}

Normalize by

\[
 y=V/V_*,\quad p=P/V_*,\quad
 \kappa=\frac{4\mu}{4-\mu},\quad b=P/(4V).
\]

Chern--Weil in the opposite orientation gives

\[
 P=\frac43V+4\int|W^-|^2-8\pi^2(2\chi-3\tau).
\]

Together with (1.1)--(1.3), this proves at every smooth time

\[
 \boxed{y\ge1,\quad y'=-p,\quad
 p\ge\Psi(y):=\max\left\{\kappa y,
                   4(1+h-y),\frac43(y-U)\right\}.}   \tag{1.5}
\]

No pointwise Ricci bound, stationarity, continuously selected monopole,
or preservation of the initial energy level is used in (1.5). The
all-metric budgets are applied afresh to each metric of the flow.

Suppose initially \(b(0)\le\bar b<1/3\). The third wall and
\(p(0)\le4\bar b\,y(0)\) imply

\[
 y(0)\le Y:=\frac U{1-3\bar b}.
\]

Thus, without any ordering assumption,

\[
                 T_{\max}\le\int_1^Y\frac{dy}{\Psi(y)}.       \tag{1.6}
\]

To see the inequality rigorously, differentiate
\(\int_1^{y(t)}\Psi(v)^{-1}dv\); its derivative is at most \(-1\). The
integral is nonnegative until the flow ceases to be smooth.

When

\[
 1\le y_L:=(1+h)(1-\mu/4)
 \le y_H:=\frac{U(4-\mu)}{4(1-\mu)}\le Y,            \tag{1.7}
\]

the maximum in (1.5) is respectively its second, first and third entry
on the three consecutive intervals. Integrating explicitly gives

\[
\boxed{\begin{aligned}
 T_{\max}\le{}&\frac14\log\frac{4h}{(1+h)\mu}\\
 &+\frac{4-\mu}{4\mu}
             \log\frac U{(1-\mu)(1+h)}\\
 &+\frac34\log\frac{4\bar b(1-\mu)}
                              {\mu(1-3\bar b)}.
\end{aligned}}                                                   \tag{1.8}
\]

These are upper bounds, not sharp lifetime claims. If (1.7) fails, use
(1.6), not the three-piece expression.

\hypertarget{a-uniform-time-bound-on-genuine-near-minimizing-metrics}{%
\subsubsection{1.3. A uniform time bound on genuine near-minimizing
metrics}\label{a-uniform-time-bound-on-genuine-near-minimizing-metrics}}

On the realized balanced family,

\[
 h_j\to0,\qquad U_j=1+h_j+O(h_j^2),\qquad
 \mathcal E(g)\ge c_1h_j\quad\hbox{for every }g,       \tag{1.9}
\]

where we may choose \(0<c_1<4\). Let scalar-minus-four initial metrics
satisfy \(b_j(0)\le\bar b_j=h_j+o(h_j)\), with \(\bar b_j\ge h_j\).
Conditions (1.7) hold for large \(j\), since

\[
 y_L=1+(1-c_1/4)h+O(h^2),\quad
 y_H=1+(1+3c_1/4)h+O(h^2),\quad
 Y=1+4h+o(h).
\]

The three terms in (1.8) tend respectively to \(\frac14\log(4/c_1)\),
\(1\), and \(\frac34\log(4/c_1)\). Therefore

\[
              \boxed{\limsup_j T_{\max,j}
                          \le1+\log\frac4{c_1}.}               \tag{1.10}
\]

The initial window is not assumed attained at the infimum. For example,

\[
 \mathcal E(g_j(0))\le\frac{4h_j}{1+h_j}+h_j^2
\]

gives \(\bar b_j=h_j+O(h_j^2)\). Such smooth scalar-minus-four metrics
exist by the inherited construction upper bound and I41's negative
Yamabe reduction. Requiring the exact upper infimum instead could give
an empty starting set when that infimum is not attained; this slack
avoids that error.

The time bound remains uniform although the all-metric defect gap tends
to zero. There is a specific reason: the two energy walls constrain the
volume to a width of order \(h\), while the linear gap forces a volume
loss rate of order \(h\). Replacing (1.9) by only a quadratic gap in the
same comparison generally gives a logarithmically divergent bound, not
(1.10).

\hypertarget{the-inherited-variational-obstruction-gives-actual-flow-descent}{%
\subsection{2. The inherited variational obstruction gives actual flow
descent}\label{the-inherited-variational-obstruction-gives-actual-flow-descent}}

This section additionally uses an actual ordinary SW solution with
projection \(p_a\ge A\) at each metric under consideration. Retain the
preceding proof's quantities

\[
 \alpha=\mathcal E=\frac{4P}{R},\quad b=\frac P{4V},\quad
 q=\frac13-2b,\quad \theta_*=-\frac1{4-\alpha}.
\]

Hold \(\theta_*\) fixed when differentiating
\(Q_{\theta_*}=R+\theta_*S_2\). For the nonpositive tensor Laplacian,
its gradient at the scalar-minus-four metric is

\[
 \mathcal G=-\Delta Z-2\operatorname{Rm}(Z)-2bZ
                                      +\tfrac12|Z|^2g,
 \qquad \operatorname{tr}\mathcal G=0.               \tag{2.1}
\]

The trace identity makes the nonlocal conformal pressure drop out of the
first variation of \(Q_{\theta_*}\). Since \(4+16\theta_*=-4b\), (1.3)
gives the exact evolution

\[
 \boxed{P'=-2\langle\mathcal G,Z\rangle_{L^2}-\frac{P^2}V
          =-2\int|\nabla Z|^2+4\int\operatorname{Rm}(Z,Z).}     \tag{2.2}
\]

Let \(C_{ijk}=\nabla_kZ_{ij}-\nabla_jZ_{ik}\), with the full ordered
tensor norm, and \(\xi=(w_3^++w_3^--1/3)_+\). The previously proved
Cotton inequality is

\[
 qP+\tfrac12\int|C|^2
              \le\langle\mathcal G,Z\rangle+\int\xi|Z|^2.
\]

Define \(D=D_2\), \(B=\int|W^-|^2\), and \(\mathcal H=D/18+3B\). The
actual same-solution estimates give \(\int\xi^2\le8\mathcal H/9\).
Inserting them into (2.2), the terms involving \(bP\) cancel exactly:

\[
 \boxed{P'\le
    \left[-\frac23+2\sqrt{\frac{8\mathcal H}{9P}}\|Z\|_\infty\right]P
                      -\int|C|^2.}                            \tag{2.3}
\]

Here and below \(\|Z\|_\infty\) is Hilbert--Schmidt, not operator norm.
Likewise direct differentiation of \(\alpha=4P/R\) gives

\[
 \boxed{-\frac{d}{dt}\log\alpha\ge\frac{4-\alpha}{2}
 \left[q-\sqrt{\frac{8\mathcal H}{9P}}\|Z\|_\infty
                                  +\frac{\int|C|^2}{2P}\right].} \tag{2.4}
\]

In the balanced family, at every smooth time,

\[
 \mathcal H=P+2\pi^2A(U-1-h-2(p_a/A-1)),
 \qquad \mathcal H/P\le1+O(h).                      \tag{2.5}
\]

The lower bound \(P\ge cAh\) follows from (1.1) and \(V\ge V_*\).
Together with \(p_a\ge A\) and \(U-1-h=O(h^2)\), it proves the ratio
bound in (2.5) without preserving an initial energy window or choosing
monopoles continuously in time. Consequently an initial sequence with
\(\|Z\|_\infty\to0\) has an actual decreasing direction along (1.2),
with initial relative rate at least \(2/3-o(1)\) in (2.4). This is
stronger than merely exhibiting a nonzero abstract gradient, but does
not assume the pinching persists.

For a finite-time formulation, consider a portion of the smooth flow on
which the following bounds hold:

\[
 b\le1/24,\quad \alpha\le1,\quad \mathcal H/P\le2,
 \quad\|Z\|_\infty\le1/32.
\]

While this tube and smoothness persist, (2.4) implies
\(-d\log\alpha/dt\ge1/4\). Starting in Section 1.3's window, the energy
level remains in that window while the estimate applies. Since
\(\alpha\ge c_1h\), the tube cannot persist beyond

\[
                     4\log\frac{\alpha(0)}{c_1h}.              \tag{2.6}
\]

Thus there is actual energy improvement until pointwise escape or loss
of smooth continuation: for sufficiently small \(h\), (2.5) prevents
exit through the \(\mathcal H/P\) face, while decreasing \(\alpha\) also
decreases \(b=\alpha/(4-\alpha)\). The unconditional bound (1.10)
remains valid after that escape. Neither conclusion presupposes a smooth
minimizer.

\hypertarget{an-actual-deformation-makes-the-weyl-payment-indispensable}{%
\subsection{3. An actual deformation makes the Weyl payment
indispensable}\label{an-actual-deformation-makes-the-weyl-payment-indispensable}}

Use the ordinary-SW quantities from the preceding proof, Section 4.1:
the scalar is \(-4\), so the spinor maximum principle gives
\(0\le\rho\le4\), hence \(0\le x\le1\).

\[
 \rho=|\Phi|^2,\quad x=\rho/4,\quad \omega=\sigma(\Phi),\quad
 K=\int|\nabla_{\mathcal A}\Phi|^2,\quad
 \gamma=\operatorname{Im}\langle\Phi,\nabla_{\mathcal A}\Phi\rangle,
\] \[
 \Gamma=\int(|d\rho|^2/4+|\gamma|^2),\quad
 J=\|F_{\mathcal A}^+-(F_{\mathcal A}^+)_H\|_2^2,
 \quad F=\int|W^+-W_0|^2,
\] \[
 W_0=\tfrac12(|\omega|^2I/3-\omega\otimes\omega),\qquad
 L_0=\tfrac23\int(1-x)^2(1+8x-2x^2).
\]

Set

\[
 \mathscr S=K-\Gamma-J-L_0-4\pi^2(p_a-A).
\]

The inherited exact signed identity is

\[
              P-8\pi^2Ah=2(2F-\mathscr S).                    \tag{3.1}
\]

One sufficient target for coefficient four had been to show
\(\mathscr S_+=o(Ah)\), discarding \(2F\). This section proves that
ordinary SW and smooth closeness to a KE metric alone cannot justify
that stronger omission. The balanced target itself is not refuted.

\hypertarget{the-closed-metric-and-its-actual-solution}{%
\subsubsection{3.1. The closed metric and its actual
solution}\label{the-closed-metric-and-its-actual-solution}}

Let \(M=\Sigma_1\times\Sigma_2\), each factor a closed hyperbolic
surface of curvature \(-1\). Its product metric \(g_0\) is negative KE,
with \(s=-4\), Kähler form \(\Omega=\Omega_1+\Omega_2\), and

\[
 A=c_1(K_M^{-1})^2=T_+=8(g_1-1)(g_2-1),\qquad h=0,\quad U=3.
\]

Its canonical ordinary SW solution has parallel spinor of squared norm
four and determinant curvature \(\pm i\Omega\). Choose a nonzero real
holomorphic quadratic differential \(h_1\) on \(\Sigma_1\), viewed as a
TT symmetric tensor, and a nonconstant real eigenfunction on
\(\Sigma_2\), \(\Delta f=\lambda f\), \(\lambda>0\). Put \(H=f h_1\),
extended by zero on other components, and \(B_0=g_0^{-1}H\). Then \(H\)
is TT on the product and \(B_0J=-JB_0\).

The smooth metrics

\[
 g_\varepsilon^{AK}(u,v)=g_0(e^{\varepsilon B_0}u,v)
\]

remain compatible with the fixed closed form \(\Omega\). Let
\(\widehat g_\varepsilon=u_\varepsilon^2g_\varepsilon^{AK}\) be their
scalar-minus-four Yamabe representatives. Negative scalar normalization
is smooth near \(g_0\). Since the scalar derivative in the TT Einstein
direction is zero, \((6\Delta+8)\dot u=0\), so

\[
 \dot u=0,\qquad \dot{\widehat g}=H.                \tag{3.2}
\]

The fixed form \(\Omega\) is harmonic and self-dual for the entire path;
self-duality of two-forms is conformally invariant in dimension four.
Thus the same canonical class has \(p_a=A\) exactly.

The canonical Kähler monopole is a regular isolated solution modulo
gauge. In the divisor description its divisor is empty; tangent and
obstruction spaces vanish. The exact classical input is
\href{https://library.slmath.org/books/Book37/files/okonek.pdf}{Okonek--Teleman,
Theorem 3.1.2, pp.~413--414}, with trivial auxiliary line and
perturbation zero. Its expected dimension is also zero. The gauge-fixed
linearization is therefore invertible. The implicit function theorem
supplies a smooth branch of genuine, unperturbed ordinary SW solutions
for \(\widehat g_\varepsilon\).

Under the natural metric identification of spinors, the metric
derivative of the Dirac equation vanishes: its three terms involve
respectively \(H\nabla_{\mathcal A_0}\Phi_0\), divergence and trace of
\(H\). The curvature equation's metric derivative also vanishes: the
trace is zero and \(F_{\mathcal A_0}^-=0\). These formulas are in
\href{https://www.numdam.org/item/10.5802/aif.2640.pdf}{Scala, Theorem
3.3 and Section 4, pp.~1274--1276}. In a fixed gauge slice,
invertibility now gives

\[
                    \dot{\mathcal A}=0,\qquad\dot\Phi=0.      \tag{3.3}
\]

There is no invocation of formal cohomological solvability or of a
perturbed monopole equation.

Regularity can also be checked directly in the project's normalization.
On the fixed KE background, write \(b_0=\mathcal A-\mathcal A_0\).
Lichnerowicz and completion of the curvature square give

\[
 \|D_{\mathcal A}\Phi\|_2^2+
 \|-iF_{\mathcal A}^+-\sigma(\Phi)\|_2^2
 =\|\nabla_{\mathcal A}\Phi\|_2^2
   +\tfrac18\||\Phi|^2-4\|_2^2+\tfrac12\|db_0\|_2^2.        \tag{3.3a}
\]

Indeed the cross terms cancel, while harmonic self-duality and exactness
give \(\|F_{\mathcal A}^+\|_2^2=2V+\frac12\|db_0\|_2^2\). A linearized
kernel element \((b_0,\phi)\) therefore satisfies

\[
 \nabla\phi+\tfrac12b_0\Phi_0=0,\quad db_0=0,\quad
 \operatorname{Re}\langle\Phi_0,\phi\rangle=0.
\]

Its component perpendicular to \(\Phi_0\) would be parallel in the
complementary line, which has nonzero curvature, and must vanish. The
remaining solution is \(\phi=it\Phi_0, b_0=-2i\,dt\), a gauge direction.
The gauge-fixed index is \((a^2-T_+)/4=0\), so the operator is
invertible. Thus the use of the implicit function theorem does not rely
on uniqueness of the nonlinear solution alone.

\hypertarget{the-losses-and-the-directional-term-have-different-orders}{%
\subsubsection{3.2. The losses and the directional term have different
orders}\label{the-losses-and-the-directional-term-have-different-orders}}

Because \(B_0\) is anti-Hermitian, its derivation action on \(\Omega\)
vanishes. Thus the derivative of the actual two-form
\(\sigma(\Phi_\varepsilon)\), as an ordinary differential form, is zero.
It follows that \(\rho-4=O(\varepsilon^2)\) and

\[
                 \Gamma,J,L_0=O(\varepsilon^4).                \tag{3.4}
\]

For \(\Gamma\), use the exact identity \(|\gamma|^2=4|d\omega|^2\); for
\(J\), use the unchanged harmonic curvature representative
\(\pm i\Omega\). All orders here hold in fixed smooth norms on this
fixed background.

The covariant derivative of \(\omega\), however, has a nonzero first
variation. At a product orthonormal frame, with \(a,b\) indices on the
first factor and \(\mu,\nu\) on the second, the mixed connection
variations are

\[
 \dot\Gamma^b_{a\mu}=\tfrac12 f_\mu(h_1)_{ab},\qquad
 \dot\Gamma^\mu_{ab}=-\tfrac12 f_\mu(h_1)_{ab}.
\]

The pure first-factor contribution to \(\dot\nabla\Omega\) vanishes:
every area-compatible two-dimensional metric preserves its area form.
The displayed mixed terms give, with the usual two-form norm,

\[
 |\dot\nabla\Omega|^2=\tfrac12|h_1|^2|df|^2.
\]

The Hopf identity
\(|\nabla\omega|^2=(\rho/2)|\nabla_{\mathcal A}\Phi|^2-|\gamma|^2/2\)
therefore proves

\[
 K=c\varepsilon^2+O(\varepsilon^3),\qquad
 c=\tfrac14\|h_1\|_2^2\|df\|_2^2>0.              \tag{3.5}
\]

Since \(h_1\) is a hyperbolic infinitesimal Einstein deformation, the
product TT variation satisfies \(\dot Z=(\lambda/2)H\). Equivalently
apply the Einstein deformation operator to \(f h_1\); only the scalar
Laplacian of \(f\) remains. Hence

\[
 \frac PK\longrightarrow\lambda,\qquad
 \frac{\mathscr S}K\longrightarrow1,\qquad
 \frac FK\longrightarrow\frac{\lambda+2}{4}.        \tag{3.6}
\]

The last limit follows as well from (3.1). The exact nonnegative
decomposition \(D/18=F+\Gamma/2+J+H_{\rm amp}\) gives

\[
             \boxed{\frac{\mathscr S_+}{P+D}
                          \longrightarrow\frac2{11\lambda+18}>0.} \tag{3.7}
\]

Thus the omission fails on genuine smooth solutions, with no zeros or
concentration needed. The missing positive contribution is nonintegrable
directional variation, paid for by \(2F\). This example has \(h=0,U=3\),
not the balanced small-h topology or its minimizing window.

An exact simplification helps isolate this distinction. With
\(T_\omega=\int(|\nabla\omega|^2-2|d\omega|^2)\), the scalar and
weighted scalar identities give

\[
 \boxed{\mathscr S=T_\omega-
       \tfrac23\int(1-x^2)(1+2x^2)-J-4\pi^2(p_a-A).}          \tag{3.8}
\]

For Kähler-aligned solutions \(T_\omega=0\), all terms on the right are
nonpositive. The actual deformation (3.5) is precisely a case where that
favorable alignment is lost.

\hypertarget{the-coefficient-of-the-weyl-payment-is-genuinely-sharp}{%
\subsubsection{3.3. The coefficient of the Weyl payment is genuinely
sharp}\label{the-coefficient-of-the-weyl-payment-is-genuinely-sharp}}

The preceding construction continues to a sharp universal statement on
one fixed smooth manifold. Fix the first surface and let the second have
genus two. Double a one-holed hyperbolic torus with geodesic boundary
length \(\ell\) along its boundary. Each half has area \(2\pi\); the
separating geodesic has the standard collar

\[
 dr^2+\ell^2\cosh^2r\,d\vartheta^2,\qquad
 \vartheta\in\mathbb R/\mathbb Z,\qquad
 |r|<\operatorname{arsinh}(1/\sinh(\ell/2)).
\]

For sufficiently small \(\ell\), this width exceeds one. The odd
Lipschitz function equal to \(r\) on \([-1,1]\) and to \(\pm1\) on the
two remaining pieces has mean zero. Direct integration gives

\[
 \int|dv|^2=2\ell\sinh1,\qquad
 \int v^2=4\pi-\ell\int_{-1}^1(1-r^2)\cosh r\,dr
          =4\pi-4\ell/e.
\]

The Rayleigh principle, valid for this \(H^1\) test, supplies actual
first eigenfunctions with

\[
 0<\lambda_1\le\frac{2\ell\sinh1}{4\pi-4\ell/e}\longrightarrow0. \tag{3.9}
\]

This uses the classical hyperbolic pants/collar construction, not a new
spectral existence theorem. The collar metric follows from Fermi
coordinates around a geodesic in curvature \(-1\).

Take \(\ell_j\to0\), normalize \(\|df_j\|_2^2=1\) and \(\|h_1\|_2^2=4\).
At every fixed \(j\), the preceding actual construction has

\[
 K_j(\varepsilon)=\varepsilon^2+O_j(\varepsilon^3),\quad
 P_j(\varepsilon)=\lambda_j\varepsilon^2+O_j(\varepsilon^3),\quad
 \Gamma_j+J_j+L_{0,j}=O_j(\varepsilon^4).
\]

Choose \(\eta_j\to0\) and then choose \(\varepsilon_j>0\) inside the
individual SW/Yamabe existence interval, with \(\varepsilon_j<1/j\), so
that all three normalized errors are at most \(\eta_j\lambda_j\):

\[
 |K_j/\varepsilon_j^2-1|,\quad
 |P_j/\varepsilon_j^2-\lambda_j|,\quad
 (\Gamma_j+J_j+L_{0,j})/\varepsilon_j^2
                                      \le\eta_j\lambda_j.
\]

No uniform inverse radius is assumed. A further decrease of
\(\varepsilon_j\) may impose arbitrarily small relative \(C^2\) distance
to that background. Since \(p_a=A\) exactly, (3.1) now proves

\[
 \frac{P_j}{\lambda_jK_j}\to1,\quad
 \frac{\mathscr S_j}{K_j}\to1,\quad
 \frac{F_j}{K_j}\to\frac12,\qquad
 \boxed{\frac{\mathscr S_j}{F_j}\to2,\quad\frac{P_j}{F_j}\to0.} \tag{3.10}
\]

At \(h=0\), (3.1) universally gives \(\mathscr S=2F-P/2\le2F\). The
sequence proves that its coefficient two cannot be reduced, even with an
arbitrary fixed finite Ricci allowance: for every \(c<2\) and finite
\(C\),

\[
                         \mathscr S\le cF+CP                   \tag{3.11}
\]

fails on sufficiently late examples. This is a sharp actual-field
statement, not a formal matrix or moment construction. It remains on
\(h=0,U=3\); the balanced coefficient-four question needs its additional
topology or minimizing geometry and remains open.

\hypertarget{the-missing-minimizing-direction-is-not-conformal}{%
\subsection{4. The missing minimizing direction is not
conformal}\label{the-missing-minimizing-direction-is-not-conformal}}

There is also an exact second-variation reduction, useful for continuing
the coefficient problem without guessing a negative conformal mode. Let
\(g\) be an actual smooth critical point of \(\mathcal E\), with

\[
 s=-4,\quad0<\alpha=\mathcal E(g)<1,\quad
 a=4-\alpha,\quad\theta=-1/a,\quad
 c=\theta+1/3=(1-\alpha)/(3a)>0.
\]

Here \(a\) is a scalar parameter, not a gauge class. In this section
\(\Delta=d^*d\) again denotes the positive scalar Laplacian. Put

\[
 L=3\Delta+4,\quad M=a\Delta+4,
 \quad\phi_0=\langle Z,k\rangle-\fint\langle Z,k\rangle
\]

for a TT tensor \(k\). For mean-zero \(f\),

\[
\boxed{\begin{aligned}
 \operatorname{Hess}\mathcal E(k+2fg,k+2fg)
 ={}&\frac aR\left[\operatorname{Hess}Q_\theta(k,k)
                   -6c\langle\phi_0,\Delta L^{-1}\phi_0\rangle\right]\\
 &+\frac8R\left\langle f-\tfrac12L^{-1}\phi_0,
                  LM(f-\tfrac12L^{-1}\phi_0)\right\rangle.
\end{aligned}}                                                   \tag{4.1}
\]

All brackets use the same metric's real \(L^2\) inner product;
\(\theta\) is held fixed in the Hessian of \(Q_\theta\).

Here is a derivation. Weyl conformal invariance gives
\(Q_\theta=2\int|W|^2+cS_2-16\pi^2\chi\). At a csc critical metric,

\[
 \operatorname{Hess}Q_\theta(2fg,2fg)=24c\langle f,\Delta Lf\rangle,
 \quad
 \operatorname{Hess}Q_\theta(k,2fg)=-12c\langle\phi_0,\Delta f\rangle.
\]

Also \(\mathcal E-J=(S_2-S_1^2/V)/R\) and \(J-\alpha=aQ_\theta/R\),
where \(J=4P/R\). Both first variations vanish at this metric. The
mean-zero scalar derivative is \((s'[k+2fg])_0=-\phi_0+2Lf\), so

\[
 \operatorname{Hess}\mathcal E(h,h)
  =\frac aR\operatorname{Hess}Q_\theta(h,h)
                  +\frac2R\|(s'[h])_0\|_2^2.
\]

Expanding and completing the square proves (4.1). In particular,

\[
 \operatorname{Hess}\mathcal E(2fg,2fg)
      =\frac8R\langle f,(3\Delta+4)(a\Delta+4)f\rangle>0       \tag{4.2}
\]

for nonzero mean-zero \(f\). Constants are the separate scaling kernel.
The minimizing conformal correction in (4.1) is exactly the one making
the first scalar variation spatially constant. Modulo diffeomorphisms
and scaling, stability therefore reduces to the actual TT form

\[
 \operatorname{Hess}Q_\theta(k,k)
            -6c\langle\phi_0,\Delta(3\Delta+4)^{-1}\phi_0\rangle. \tag{4.3}
\]

The correction is bounded and nonnegative before its minus sign:
\(0\le\Delta(3\Delta+4)^{-1}<1/3\). This identifies the remaining
nonlinear variational question; it does not prove (4.3) is negative. No
Hessian formula valid only at an Einstein background is substituted at
the non-Einstein candidates.

\hypertarget{a-stationary-zero-end-sequence-must-lose-sobolev-control}{%
\subsection{5. A stationary zero-end sequence must lose Sobolev
control}\label{a-stationary-zero-end-sequence-must-lose-sobolev-control}}

This is an alternative elliptic endpoint result, not a deduction of
I48's small local lower-Ricci budget. Fix \(S,L_0,\Theta,T_0,\chi_0\).
Let

\[
 s=-4,\quad \nabla Q_\theta=0,\quad |\theta|\le\Theta,
 \quad 0<T_+=2\chi+3\tau\le T_0,\quad\chi\le\chi_0,          \tag{5.1}
\]

on closed connected oriented four-manifolds, and assume the genuine
global Sobolev inequality

\[
 \left(\int|u|^4\right)^{1/2}
       \le S\int|du|^2+V^{-1/2}\int|u|^2.                     \tag{5.2}
\]

For a smooth self-dual form with \(|\omega|\le L_0\), define

\[
 Z_t=|\operatorname{Ric}+g|^2+|W^+-A_{\omega,t}|^2
                    +(|\omega|^2-2)^2+|d\omega|^2,
 \quad A_{\omega,t}=t(|\omega|^2I/3-\omega\otimes\omega).
\]

There is \(c=c(S,L_0,\Theta,T_0,\chi_0)>0\) such that

\[
                    \fint Z_t\ge c,\qquad0\le t\le1/4.       \tag{5.3}
\]

\textbf{Proof.} Put \(z=\fint Z_t\). If \(z\le1/3\), then \(P\le zV\),
\(W_+\le2zV+L_0^4V/12\), and Chern--Weil gives

\[
 4\pi^2T_+=2W_++\tfrac23V-\tfrac12P,
 \qquad
 \frac{4\pi^2}{2+L_0^4/6}\le V\le8\pi^2T_0.                  \tag{5.4}
\]

The lower bound uses the positive integer \(T_+\ge1\). In the full
four-index norm,

\[
 \int|\operatorname{Rm}|^2=32\pi^2\chi+4P
                       \le32\pi^2\chi_0+\tfrac43(8\pi^2T_0). \tag{5.5}
\]

Thus normalized vanishing error implies absolute vanishing error, but
only after (5.4). At constant scalar curvature the critical equation is

\[
 \Delta\operatorname{Ric}=-2\operatorname{Rm}(\operatorname{Ric})
 -2\theta s\operatorname{Ric}
            +\tfrac12(|\operatorname{Ric}|^2+\theta s^2)g,     \tag{5.6}
\]

with nonpositive tensor Laplacian. Every right-hand term is an algebraic
\(\operatorname{Rm}*\operatorname{Ric}\) term, uniformly for bounded
\(\theta\). This is \href{https://arxiv.org/pdf/1404.0569}{Catino,
equation (2.1)}, not an Einstein equation assigned to the original
metrics.

Normalize volume to one. Equations (5.2), (5.4), (5.5) and (5.6) give
orbifold compactness and smooth convergence on regular charts by
\href{https://arxiv.org/pdf/math/0312111}{Anderson, Theorem 1.1 and
Proposition 3.5}. The Sobolev inequality supplies the lower ball-volume
bound; the elliptic equation supplies the small-energy estimate. There
is no division by \(\theta+1/3\), and csc is retained even at that
parameter. Pass to a subsequence with \(V_j\to V_\infty\in(0,\infty)\),
justified by (5.4). Returning from \(\widehat g_j=V_j^{-1/2}g_j\) to
physical scale, a sequence violating (5.3) has a nonempty smooth regular
chart with \(\operatorname{Ric}=-g\) in the limit.

The local first-order elliptic estimate for self-dual forms bounds
\(\omega_j\) in \(H^1\), since \(d^*\omega_j=-*d\omega_j\). Rellich
compactness and the uniform amplitude bound give strong \(L^2\)
convergence of their quadratic models. The limiting form is harmonic,
smooth, and satisfies

\[
 d\omega=0,\quad |\omega|^2=2,\quad
 W^+=t_\infty(2I/3-\omega\otimes\omega),
 \quad t_\infty\in[0,1/4].
\]

The local incompatibility already proved in I42/I45 excludes this: at
\(t=0\) use the local ASD negative-Einstein almost-Kähler obstruction;
at \(t>0\), Einstein Bianchi and the constant two-eigenvalue model force
the parallel Kähler branch \(t=1/2\). This proves (5.3). QED.

On a fixed topology, an actual csc quadratic-critical zero-end sequence
must therefore lose (5.2). This does not say that an arbitrary
minimizing sequence is critical, does not remove all compactness
hypotheses, and does not address the \(\beta=2\) coefficient problem.
The topology bound is on \(T_+\) and \(\chi\) themselves, not merely
their ratio.

\hypertarget{what-curvature-maximal-extraction-supplies-and-what-it-loses}{%
\subsection{6. What curvature-maximal extraction supplies, and what it
loses}\label{what-curvature-maximal-extraction-supplies-and-what-it-loses}}

A separate bounded statement advances the core question without assuming
the desired core in advance. Let \(g_j\) be exact scalar-minus-four
critical metrics of \(Q_{\theta_j}\), with bounded \(\theta_j\),

\[
 \operatorname{Vol}B_s(x)\ge v s^4\ (0<s\le r_0),\qquad
 \int|\operatorname{Rm}|^2\le E_0,                            \tag{6.1}
\]

and bounded forms \(|\omega_j|\le L_0\), \(t_j\in[0,1/4]\), such that

\[
 \int\bigl(|\operatorname{Ric}+g|^2+
                      |W^+-A_{\omega_j,t_j}|^2\bigr)\to0.     \tag{6.2}
\]

The errors in (6.2) are absolute, not volume averages. If
\(Q_j=\max|\operatorname{Rm}_{g_j}|\to\infty\), choose maximum points
and rescale \(h_j=r_j^{-2}g_j\), \(r_j=Q_j^{-1/2}\). Then
\(|\operatorname{Rm}_{h_j}|\le1\), with equality at the base point. The
rescaled balls are noncollapsed at every fixed radius. Bounded curvature
and noncollapse give injectivity control. Harmonic coordinates and (5.6)
then bootstrap to smooth pointed convergence on compact sets to a
complete nonflat manifold \(X\).

On fixed rescaled balls, the original volume is \(O(r_j^4)\). Curvature
square energy is scale invariant; the bounded backgrounds in (6.2)
disappear. Therefore \(X\) is Ricci-flat and \(W_X^+=0\). It has finite
curvature energy and Euclidean lower volume growth. The classical
\href{https://doi.org/10.1007/BF01389045}{Bando--Kasue--Nakajima ALE
theorem} now applies to this exact Ricci-flat limit, not to \(g_j\). The
limit has one end: multiple ends would force an isometric splitting with
a compact transverse factor and hence linear volume growth, contrary to
its Euclidean lower growth. Thus the truncations below have one outer
boundary, as required by the entire-core theorem.

Large compact ALE truncations embed into \(M_j\). A diagonal choice of
truncation radius gives actual entire smooth cores whose original
volumes tend to zero and whose lifted outer collars converge to a flat
cone. Their total

\[
 \int\bigl(|\operatorname{Ric}^{\circ}|^2+4|W^+|^2+s^2/12\bigr)
                                                               \tag{6.3}
\]

tends to zero by (6.2). Explicitly its integral over a core \(\Omega_j\)
is at most
\(8\,\mathrm{error}_j+C(L_0)\operatorname{Vol}_{g_j}(\Omega_j)\). No
bubble-tree energy identity is being assumed. I47's entire-core theorem
therefore excludes the physical quotients \(1/n(1,-2)\), odd \(n\ge3\),
as their infinity groups. The group cannot be trivial either: a complete
Ricci-flat ALE space with Euclidean asymptotic volume ratio is Euclidean
by volume rigidity. Its order is bounded by \(\omega_4/v\).

The metric core has genuinely been extracted. The nonzero actual
opposite-spinor field has not. Its amplitude and Dirac errors have
different rescaling weights: absolute errors tending to zero need not
tend to zero after division by powers of \(r_j\). For instance, the mean
amplitude error on a fixed rescaled ball contains

\[
                      r_j^{-4}\int_{\Omega_j}(\rho-\rho_*)^2.
\]

Consequently it is not legitimate to combine the extracted core with
I46's simultaneous nonzero-spinor holonomy conclusion without a new
same-scale field estimate. These are terminal cores, not an identified
entire exceptional/handle surgery region or a proved bubble tree.

\hypertarget{the-next-exact-questions-left-by-these-arguments}{%
\subsection{7. The next exact questions left by these
arguments}\label{the-next-exact-questions-left-by-these-arguments}}

For the balanced coefficient, retain the full signed target
\((\mathscr S-2F)_+=o(Ah)\) in (3.1), or find a variational/core
argument producing the same conclusion. The stronger omission of \(2F\)
requires topology-specific information: it is false as a universal
ordinary-SW principle. A negative direction for a hypothetical minimizer
must be found in the reduced TT form (4.3), not in conformal directions.

For the general integral zero endpoint, a metric core and a calibrated
field on that core are distinct outputs. Section 6 makes their scaling
gap explicit. Section 5 excludes uniformly Sobolev-controlled stationary
escape, while I48 treats a different integral-Ricci class. Neither turns
raw weak input into a complete core decomposition.

For the dynamic connection, (1.10) already applies to actual smooth
near-minimizing initial metrics without a minimizer. Determining the
forced loss of smooth continuation, or proving persistence estimates
strong enough to contradict it in a broader input class, requires new
analysis. It is not inferred from the monotonic volume formula alone.

END COMPLETE SOURCE

\end{document}
